
\documentclass[a4paper,11pt]{article}
\usepackage[a4paper, margin=8em]{geometry}

% usa i pacchetti per la scrittura in italiano
\usepackage[french,italian]{babel}
\usepackage[T1]{fontenc}
\usepackage[utf8]{inputenc}
\frenchspacing 

% usa i pacchetti per la formattazione matematica
\usepackage{amsmath, amssymb, amsthm, amsfonts}

% usa altri pacchetti
\usepackage{gensymb}
\usepackage{hyperref}
\usepackage{standalone}

\usepackage{colortbl}

\usepackage{xstring}
\usepackage{karnaugh-map}

% imposta il titolo
\title{Embedded Operating Systems and Architectures notes}
\author{Luca Seggiani}
\date{2026}

% imposta lo stile
% usa helvetica
\usepackage[scaled]{helvet}
% usa palatino
\usepackage{palatino}
% usa un font monospazio guardabile
\usepackage{lmodern}

\renewcommand{\rmdefault}{ppl}
\renewcommand{\sfdefault}{phv}
\renewcommand{\ttdefault}{lmtt}

% circuiti
\usepackage{circuitikz}
\usetikzlibrary{babel}

% testo cerchiato
\newcommand*\circled[1]{\tikz[baseline=(char.base)]{
            \node[shape=circle,draw,inner sep=2pt] (char) {#1};}}

% disponi il titolo
\makeatletter
\renewcommand{\maketitle} {
	\begin{center} 
		\begin{minipage}[t]{.8\textwidth}
			\textsf{\huge\bfseries \@title} 
		\end{minipage}%
		\begin{minipage}[t]{.2\textwidth}
			\raggedleft \vspace{-1.65em}
			\textsf{\small \@author} \vfill
			\textsf{\small \@date}
		\end{minipage}
		\par
	\end{center}

	\thispagestyle{empty}
	\pagestyle{fancy}
}
\makeatother

% disponi teoremi
\usepackage{tcolorbox}
\newtcolorbox[auto counter, number within=section]{theorem}[2][]{%
	colback=blue!10, 
	colframe=blue!40!black, 
	sharp corners=northwest,
	fonttitle=\sffamily\bfseries, 
	title=Theorem~\thetcbcounter: #2, 
	#1
}

% disponi definizioni
\newtcolorbox[auto counter, number within=section]{definition}[2][]{%
	colback=red!10,
	colframe=red!40!black,
	sharp corners=northwest,
	fonttitle=\sffamily\bfseries,
	title=Definition~\thetcbcounter: #2,
	#1
}

% disponi codice
\usepackage{listings}
\usepackage[table]{xcolor}

\definecolor{codegreen}{rgb}{0,0.6,0}
\definecolor{codegray}{rgb}{0.5,0.5,0.5}
\definecolor{codepurple}{rgb}{0.58,0,0.82}
\definecolor{backcolour}{rgb}{0.95,0.95,0.92}

\lstdefinestyle{codestyle}{
		backgroundcolor=\color{black!5}, 
		commentstyle=\color{codegreen},
		keywordstyle=\bfseries\color{magenta},
		numberstyle=\sffamily\tiny\color{black!60},
		stringstyle=\color{green!50!black},
		basicstyle=\ttfamily\footnotesize,
		breakatwhitespace=false,         
		breaklines=true,                 
		captionpos=b,                    
		keepspaces=true,                 
		numbers=left,                    
		numbersep=5pt,                  
		showspaces=false,                
		showstringspaces=false,
		showtabs=false,                  
		tabsize=2
}

\lstdefinestyle{shellstyle}{
		backgroundcolor=\color{black!5}, 
		basicstyle=\ttfamily\footnotesize\color{black}, 
		commentstyle=\color{black}, 
		keywordstyle=\color{black},
		numberstyle=\color{black!5},
		stringstyle=\color{black}, 
		showspaces=false,
		showstringspaces=false, 
		showtabs=false, 
		tabsize=2, 
		numbers=none, 
		breaklines=true
}

\lstdefinelanguage{x86}{ 
  keywords={AAA, AAD, AAM, AAS, ADC, ADCB, ADCW, ADCL, ADD, ADDB, ADDW, ADDL, AND, ANDB, ANDW, ANDL,
        ARPL, BOUND, BSF, BSFL, BSFW, BSR, BSRL, BSRW, BSWAP, BT, BTC, BTCB, BTCW, BTCL, BTR, 
        BTRB, BTRW, BTRL, BTS, BTSB, BTSW, BTSL, CALL, CBW, CDQ, CLC, CLD, CLI, CLTS, CMC, CMP,
        CMPB, CMPW, CMPL, CMPS, CMPSB, CMPSD, CMPSW, CMPXCHG, CMPXCHGB, CMPXCHGW, CMPXCHGL,
        CMPXCHG8B, CPUID, CWDE, DAA, DAS, DEC, DECB, DECW, DECL, DIV, DIVB, DIVW, DIVL, ENTER,
        HLT, IDIV, IDIVB, IDIVW, IDIVL, IMUL, IMULB, IMULW, IMULL, IN, INB, INW, INL, INC, INCB,
        INCW, INCL, INS, INSB, INSD, INSW, INT, INT3, INTO, INVD, INVLPG, IRET, IRETD, JA, JAE,
        JB, JBE, JC, JCXZ, JE, JECXZ, JG, JGE, JL, JLE, JMP, JNA, JNAE, JNB, JNBE, JNC, JNE, JNG,
        JNGE, JNL, JNLE, JNO, JNP, JNS, JNZ, JO, JP, JPE, JPO, JS, JZ, LAHF, LAR, LCALL, LDS,
        LEA, LEAVE, LES, LFS, LGDT, LGS, LIDT, LMSW, LOCK, LODSB, LODSD, LODSW, LOOP, LOOPE,
        LOOPNE, LSL, LSS, LTR, MOV, MOVB, MOVW, MOVL, MOVSB, MOVSD, MOVSW, MOVSX, MOVSXB,
        MOVSXW, MOVSXL, MOVZX, MOVZXB, MOVZXW, MOVZXL, MUL, MULB, MULW, MULL, NEG, NEGB, NEGW,
        NEGL, NOP, NOT, NOTB, NOTW, NOTL, OR, ORB, ORW, ORL, OUT, OUTB, OUTW, OUTL, OUTSB, OUTSD,
        OUTSW, POP, POPL, POPW, POPB, POPA, POPAD, POPF, POPFD, PUSH, PUSHL, PUSHW, PUSHB, PUSHA, 
        RCLL, RCR, RCRB, RCRW, RCRL, RDMSR, RDPMC, RDTSC, REP, REPE, REPNE, RET, ROL, ROLB, ROLW,
        ROLL, ROR, RORB, RORW, RORL, SAHF, SAL, SALB, SALW, SALL, SAR, SARB, SARW, SARL, SBB,
        SBBB, SBBW, SBBL, SCASB, SCASD, SCASW, SETA, SETAE, SETB, SETBE, SETC, SETE, SETG, SETGE,
        SETL, SETLE, SETNA, SETNAE, SETNB, SETNBE, SETNC, SETNE, SETNG, SETNGE, SETNL, SETNLE,
        SETNO, SETNP, SETNS, SETNZ, SETO, SETP, SETPE, SETPO, SETS, SETZ, SGDT, SHL, SHLB, SHLW,
        SHLL, SHLD, SHR, SHRB, SHRW, SHRL, SHRD, SIDT, SLDT, SMSW, STC, STD, STI, STOSB, STOSD,
        STOSW, STR, SUB, SUBB, SUBW, SUBL, TEST, TESTB, TESTW, TESTL, VERR, VERW, WAIT, WBINVD,
        XADD, XADDB, XADDW, XADDL, XCHG, XCHGB, XCHGW, XCHGL, XLAT, XLATB, XOR, XORB, XORW, XORL},
  keywordstyle=\color{blue}\bfseries,
  ndkeywordstyle=\color{darkgray}\bfseries,
  identifierstyle=\color{black},
  sensitive=false,
  comment=[l]{\#},
  morecomment=[s]{/*}{*/},
  commentstyle=\color{purple}\ttfamily,
  stringstyle=\color{red}\ttfamily,
  morestring=[b]',
  morestring=[b]"
}

\lstdefinelanguage{riscv}{
  keywords={
    LUI, AUIPC,
    JAL, JALR,
    BEQ, BNE, BLT, BGE, BLTU, BGEU,
    LB, LH, LW, LBU, LHU, LD,
    SB, SH, SW, SD,
    ADDI, SLTI, SLTIU, XORI, ORI, ANDI,
    SLLI, SRLI, SRAI,
    ADD, SUB, SLL, SLT, SLTU, XOR, SRL, SRA, OR, AND,
    ADDIW, SLLIW, SRLIW, SRAIW,
    ADDW, SUBW, SLLW, SRLW, SRAW,
    MUL, MULH, MULHSU, MULHU, DIV, DIVU, REM, REMU,
    MULW, DIVW, DIVUW, REMW, REMUW,
    LR.W, SC.W, AMOSWAP.W, AMOADD.W, AMOXOR.W,
    AMOAND.W, AMOOR.W, AMOMIN.W, AMOMAX.W,
    AMOMINU.W, AMOMAXU.W,
    LR.D, SC.D, AMOSWAP.D, AMOADD.D, AMOXOR.D,
    AMOAND.D, AMOOR.D, AMOMIN.D, AMOMAX.D,
    AMOMINU.D, AMOMAXU.D,
    FLW, FSW,
    FMADD.S, FMSUB.S, FNMSUB.S, FNMADD.S,
    FADD.S, FSUB.S, FMUL.S, FDIV.S, FSQRT.S,
    FSGNJ.S, FSGNJN.S, FSGNJX.S,
    FMIN.S, FMAX.S,
    FCVT.W.S, FCVT.WU.S, FCVT.S.W, FCVT.S.WU,
    FMV.X.W, FMV.W.X,
    FEQ.S, FLT.S, FLE.S,
    FLD, FSD,
    FMADD.D, FMSUB.D, FNMSUB.D, FNMADD.D,
    FADD.D, FSUB.D, FMUL.D, FDIV.D, FSQRT.D,
    FSGNJ.D, FSGNJN.D, FSGNJX.D,
    FMIN.D, FMAX.D,
    FCVT.W.D, FCVT.WU.D, FCVT.D.W, FCVT.D.WU,
    FCVT.S.D, FCVT.D.S,
    FMV.X.D, FMV.D.X,
    FEQ.D, FLT.D, FLE.D,
    CSRRW, CSRRS, CSRRC,
    CSRRWI, CSRRSI, CSRRCI,
    FENCE.I,
    FENCE,
    ECALL, EBREAK,
    MRET, SRET, URET,
    WFI,
    SFENCE.VMA,
    NOP, LI, LA, MV,
    NOT, NEG, NEGW,
    SEQZ, SNEZ, SLTZ, SGTZ,
    BEQZ, BNEZ, BLEZ, BGEZ, BLTZ, BGTZ,
    BGT, BLE, BGTU, BLEU,
    J, JR, RET,
    CALL, TAIL,
    RDINSTRET, RDCYCLE, RDTIME,
    CSRR, CSRW, CSRS, CSRC,
    CSRWI, CSRSI, CSRCI
  },
  keywordstyle=\color{blue}\bfseries,
  ndkeywordstyle=\color{darkgray}\bfseries,
  identifierstyle=\color{black},
  sensitive=false,
  comment=[l]{\#},
  morecomment=[s]{/*}{*/},
  commentstyle=\color{purple}\ttfamily,
  stringstyle=\color{red}\ttfamily,
  morestring=[b]',
  morestring=[b]"
}

\lstdefinelanguage{arm}{
  keywords={
    ADC, ADD, AND, BIC, CMN, CMP, EOR, MOV, MVN,
    RSB, RSC, SBC, SUB, TST, TEQ,
    MLA, MLS, MUL, SMLAL, SMULL, UMLAL, UMULL,
    QADD, QSUB, QDADD, QDSUB,
    QADD16, QADD8, QASX, QSAX,
    QSUB16, QSUB8, QSAX, QSUBADDX,
    ASR, LSL, LSR, ROR, RRX,
    B, BL, BX, BLX,
    BXJ,
    BEQ, BNE, BCS, BHS, BCC, BLO,
    BMI, BPL, BVS, BVC,
    BHI, BLS, BGE, BLT, BGT, BLE,
    BAL,
    LDR, LDRB, LDRH, LDRSB, LDRSH,
    LDRD, LDM, LDMIA, LDMIB, LDMDA, LDMDB,
    LDRT, LDRBT, LDRHT, LDRSBT, LDRSHT,
    STR, STRB, STRH, STRD,
    STM, STMIA, STMIB, STMDA, STMDB,
    STRT, STRBT, STRHT,
    PUSH, POP,
    SWP, SWPB,
    REV, REV16, REVSH,
    BFC, BFI, BFXIL, SBFX, UBFX,
    CLZ,
    LDREX, LDREXB, LDREXH, LDREXD,
    STREX, STREXB, STREXH, STREXD,
    DMB, DSB, ISB,
    MCR, MCRR, MRC, MRRC,
    MRS, MSR,
    CPS, SETEND, SVC, BKPT,
    CLREX, WFE, WFI, SEV,
    YIELD, NOP,
    VLDR, VSTR,
    VMOV, VABS, VNEG,
    VADD, VSUB, VMUL, VDIV,
    VMLA, VMLS, VNMLA, VNMLS, VNMUL,
    VMAX, VMIN,
    VSQRT,
    VCMP, VCMPE,
    VCVT,
    VCLZ, VCLS,
    VAND, VORR, VEOR, VBIC, VMVN,
    VLD1, VLD2, VLD3, VLD4,
    VST1, VST2, VST3, VST4,
    ADR, ADCS, ADDS, SUBS, MOVS, ANDS, ORRS, EORS
  },
  keywordstyle=\color{blue}\bfseries,
  ndkeywordstyle=\color{darkgray}\bfseries,
  identifierstyle=\color{black},
  sensitive=false,
  comment=[l]{@},
  morecomment=[l]{;},
  morecomment=[s]{/*}{*/},
  commentstyle=\color{purple}\ttfamily,
  stringstyle=\color{red}\ttfamily,
  morestring=[b]',
  morestring=[b]"
}

\lstset{style=codestyle, language=C}

% disponi sezioni
\usepackage{titlesec}

\titleformat{\section}
	{\sffamily\Large\bfseries} 
	{\thesection}{1em}{} 
\titleformat{\subsection}
	{\sffamily\large\bfseries}   
	{\thesubsection}{1em}{} 
\titleformat{\subsubsection}
	{\sffamily\normalsize\bfseries} 
	{\thesubsubsection}{1em}{}

% tikz
\usepackage{tikz}

% float
\usepackage{float}

% grafici
\usepackage{pgfplots}
\pgfplotsset{width=10cm,compat=1.9}

% disponi alberi
\usepackage{forest}

\forestset{
	rectstyle/.style={
		for tree={rectangle,draw,font=\large\sffamily}
	},
	roundstyle/.style={
		for tree={circle,draw,font=\large}
	}
}

% disponi algoritmi
\usepackage{algorithm}
\usepackage{algorithmic}
\makeatletter
\renewcommand{\ALG@name}{Algorithm}
\makeatother

% disponi numeri di pagina
\usepackage{fancyhdr}
\fancyhf{} 
\fancyfoot[L]{\sffamily{\thepage}}

\makeatletter
\fancyhead[L]{\raisebox{1ex}[0pt][0pt]{\sffamily{\@title \ \@date}}} 
\fancyhead[R]{\raisebox{1ex}[0pt][0pt]{\sffamily{\@author}}}
\makeatother

\begin{document}
% sezione (data)
\section{Lesson 23-09-26}

% stili pagina
\thispagestyle{empty}
\pagestyle{fancy}

% testo
\subsection{Introduction}

Computing systems are everywhere: cars, trains, aircraft, satellites, medical devices, printers, phones, industrial controllers.
In an \textbf{embedded} system, we consider \textit{software plus hardware}, not software alone: the program only makes sense together with the machine it drives.
Complexity is growing dramatically.
A modern car carries $\approx$ 100 electronic control units and on the order of $10^8$ lines of on-board code, against a handful of hundreds of lines in the 1970s.

\subsubsection{Operating systems}

The natural question to ask is \textit{why} an embedded system might need an operating system.
The answers are many:

\begin{itemize}
	\item
	      Bare metal applications no longer cope with the complexity even in
	      the embedded systems domain;
	\item
	      \textbf{An operating system raises the abstraction level:} it gives
	      application code a small set of rules and services instead of ad-hoc
	      hardware handling;
\end{itemize}
Our goal is to connect three things that are usually taught apart:
\textbf{computer architecture}, \textbf{operating systems}, and the
\textbf{computer system} that has to run both.

\subsubsection{Topics}

The main roadmap of the course is:

\begin{itemize}
	\item
	      We will connect operating-system concepts to executable embedded
	      software;
	\item
	      This course lives in the \textbf{platform} layer of a computing node:
	      from the hardware up to the system-call interface, with the
	      application used as the thing that exercises it.
\end{itemize}

The covered topics are:

\begin{itemize}
	\item
	      \textbf{Minimal computing-system architecture} --- juste enough to
	      write operating systems code;
	\item
	      \textbf{Bare-metal development flow} --- cross-compilation, emulators
	      and simulators, debugging and profiling, virtual platforms;
	\item
	      \textbf{Operating systems for embedded systems} --- real-time and
	      process scheduling, the anatomy of an embedded OS, architecture and
	      services of real systems (FreeRTOS);
	\item
	      \textbf{Programming an embedded OS} --- building it, device management
	      and drivers, drivers for custom devices.
\end{itemize}

\subsection{Introduction to embedded operating systems}

An \textbf{OS} (\textit{Operating System}) is nothing but a program, with a very simple purpose: to mediate between hardware and the programs that use it.

There are three classic ways to draw that boundary:
\begin{itemize}
	\item \textbf{Monolithic}: kernel is one large program; drivers, filesystems, scheduling all share one address space. Both Linux and Windows fit this description.

	      Monolithic kernels are simple to implement once the task switching machinery is running: once we can switch to kernel mode, we can touch any resource available to the system to carry out tasks;
	\item \textbf{Microkernel}: kernel does the bare minimum (IPC, scheduling, basic memory); everything else runs as separate services.

	      Microkernels are interesting because they are inherently safer: the service handling the filesystem, for instance, cannot corrupt other services on failure (as it runs itself as a protected service with no access to other resources).
	      The problem comes up when complex tasks have to performed: the main overhead is given by the IPC and scheduling facilities which sit inbetween services;
	\item \textbf{Flat}: no kernel/user separation at all; a small set of statically-defined tasks share the CPU under one scheduler.

	      This is important to us because embedded systems often don't feature protection systems, and as we know complete protection is impossible without hardware support.
\end{itemize}

The core of the operating system is basically the \textit{scheduler}: everything else can be seen as a task either the operating system or the user programs carry out, based on the resources allocated by the scheduler.

\subsubsection{Hardware resources}

The purpose of an OS is to make hardware \textit{usable}.
Without an OS, every application must decide how to initialise the CPU, program a timer, share a UART, and recover from an interrupt.
An OS supplies a small set of rules and mechanisms so application code can ask for useful services instead of repeatedly rebuilding hardware policy.

The resources are finite, especially on a small board:
\begin{itemize}
	\item \textbf{CPU time}: which ready task runs now? We also haveto take into account the fact that CPUs interact with many other devices on the same board;
	\item \textbf{memory}: which code and data occupy RAM and Flash?
	\item \textbf{devices}: who may configure a timer or consume a received byte?
	\item \textbf{events}: how does an interrupt become work for the right task?
\end{itemize}

\subsubsection{Types of operating systems}

The three labels we gave earlier describe where operating-system services execute and
how strongly they are isolated.

\begin{table}[H]
	\center \rowcolors{2}{white}{black!10}
	\begin{tabular} { c | c | c }
		\bfseries Type & \bfseries Services offered     & \bfseries Tradeoff          \\
		\hline
		Monolithic     & Together in kernel space       & Direct calls, large trusted
		core                                                                          \\
		Microkernel    & Mostly as separate servers     & Isolation, IPC
		overhead/design                                                               \\
		Flat           & Together in one firmware image & Tiny and direct, little
		isolation                                                                     \\
	\end{tabular}
\end{table}

\textbf{\textsf{Monolithic}}
A monolithic kernel puts scheduling, drivers, filesystems, and
networking in one kernel address space. Services can call each other
directly and quickly.

\begin{itemize}
	\item
	      Linux is the familiar large-scale example
	\item
	      A faulty driver can corrupt kernel state
	\item
	      The design is practical when the feature set is broad and tightly
	      coupled
\end{itemize}

On a protected machine, applications normally execute outside this
trusted kernel address space.

\par\smallskip

\textbf{\textsf{Microkernel}}
A microkernel retains only the most fundamental mechanisms in the
kernel, usually scheduling, address-space management, and inter-process
communication.

Drivers and other services become separate programs that exchange
messages.

\begin{itemize}
	\item
	      Smaller trusted computing base
	\item
	      A service failure can be contained or restarted
	\item
	      More interfaces and message paths to design correctly
\end{itemize}

seL4 and QNX are useful reference points for this family.

\par\smallskip

\textbf{\textsf{Flat}}
Flat embedded firmware usually has one linked image and one physical
address space. Tasks cooperate through shared memory and the scheduler.

There is no process loader, no virtual-memory translation, and usually
no hardware-enforced boundary between tasks.
This puts more pressure on software developers, as user code must be free of bugs which could potentially crash the whole operating system.
This is particularly critical in systems like vehicles, medical or industrial.
We know that monolithic operating systems give us an extra layer of safety, by switching to a safer kernel mode to perform critical tasks, which user code can't normally access.

\subsubsection{Operating system roadmap}

Where does this course sit in this hierachy?

\begin{itemize}
	\item
	      \textbf{Linux} (a monolithic kernel) and \textit{seL4/QNX} (microkernels) are the extremes we'll read about;
	\item
	      This course builds \textbf{flat} first (\texttt{scratch-os}) --- the simplest point on the spectrum, and the one that makes the boot-to-scheduling pipeline fully visible.
	      \texttt{scratch-os} is intentionally small. Its value is that every important step
	      can fit in mind;
	\item
	      Then \textbf{FreeRTOS} --- closer to a minimal microkernel-flavored RTOS: a small kernel, real scheduling policies, but still no memory protection between tasks by default.

	      FreeRTOS provides established implementations of the mechanisms we first
	      meet in \texttt{scratch-os}:

	      \begin{itemize}
		      \item
		            Task creation and priorities;
		      \item
		            Delay and timer-driven scheduling;
		      \item
		            Queues, semaphores, and mutexes;
		      \item
		            Interrupt-to-task communication.
	      \end{itemize}

	      The comparison matters: an API call becomes easier to reason about when
	      its underlying state transitions are no longer invisible.
\end{itemize}

\subsubsection{Real-time operating systems}

The \textbf{real-time} label applied to an OS is about \textit{bounded response}, not raw speed.
The idea is that we can know exactly how long a task will take before executing it.
A result computed after its deadline can be a failure even when its value is otherwise correct.
This is particularly useful in embedded systems which might have external (e.g. physical) constranints, in addition to technical ones.

\begin{table}[h!]
	\center \rowcolors{2}{white}{black!10}
	\begin{tabular} { c | c  }
		\bfseries Requirement & \bfseries Example                         \\
		\hline
		Functional            & sample a sensor correctly                 \\
		Temporal              & sample it every 1 ms, before the deadline \\
		Robustness            & continue safely when one input is missing \\
	\end{tabular}
\end{table}

\subsubsection{Embedded operating systems}

\textbf{Embedded} operating systems usually look different than traditional ones:

\begin{itemize}
	\item
	      \textbf{Resource-constrained}: kilobytes of RAM/flash, not gigabytes
	      --- no room for a general-purpose kernel's overhead;
	\item
	      \textbf{Real-time}: correctness often depends on \emph{when} something
	      happens, not just \emph{whether} it happens;
	\item
	      \textbf{No user protection needed (usually)}: one vendor, one firmware
	      image, one set of tasks --- the isolation problems a desktop OS solves
	      mostly don't apply here.

	      This is mainly a cost problem: it's not that car manufactures don't \textit{want} safety in their devices, but that that CPUs which don't feature protection mechanisms are usually cheaper;
	\item
	      These constraints are why \texttt{scratch-os} can be simple and still
	      be a real, working OS.
\end{itemize}

An embedded OS is selected and configured around a particular product,
not an unknown future workload.

\begin{itemize}
	\item
	      Static task counts and fixed buffers can replace dynamic allocation;
	\item
	      A simple priority scheduler can replace a general-purpose fairness
	      policy;
	\item
	      Direct register access can replace a layered device stack;
	\item
	      Predictable execution is often more valuable than maximum throughput.
\end{itemize}

Small is not incomplete when it matches the system's responsibilities.

Many \textbf{MCU} (\textit{Micro Controller Unit}) products ship one firmware image that includes startup code,
drivers, the OS, and application logic. Updating that image changes the
whole system.

This simplifies deployment, but it raises the bar for discipline:

\begin{itemize}
	\item
	      a bad pointer can damage any writable state
	\item
	      a blocked high-priority task can stop an important service
	\item
	      recovery and observability must be designed into the firmware
\end{itemize}

Later, we will see how scheduling and synchronisation reduce these
risks.

\subsubsection{Embedded systems}

There are many ways we can classify \textbf{embedded systems}:

\begin{itemize}
	\item
	      A computer built into a larger device to do \textbf{one dedicated job}
	      --- not a general-purpose machine you install arbitrary software on;
	\item
	      Hardware + software (often called \emph{firmware}) shipped as one
	      unit; the end user never sees a shell or a package manager;
	\item
	      Everywhere: engine controllers, pacemakers, drones, washing machines,
	      Wi-Fi routers, the microcontroller in a USB charger;
	\item
	      The population of embedded processors passed the human population
	      decades ago --- most CPUs made are never in a ``computer''.
\end{itemize}

We can compare embedded systems to general-purpose ones (such as modern computer or phone operating systems):

\begin{table}[h!]
	\center \rowcolors{2}{white}{black!10}
	\begin{tabular} { c | c | c }
		           & \bfseries Embedded   & \bfseries General-purpose  \\
		\hline
		Workload   & one fixed task       & anything the user installs \\
		Resources  & KB of RAM/flash      & GB of RAM, TB of disk      \\
		Timing     & often hard deadlines & best-effort                \\
		User model & no interactive user  & multi-user, multi-tasking  \\
		Cost/power & pennies, milliwatts  & dollars, tens of watts     \\
	\end{tabular}
\end{table}

The right-hand column is what a desktop OS is built for;
\texttt{scratch-os} targets the left-hand one, which is \emph{why} it
can be small.

\subsection{Microcontrollers}

We begin to approach the device that we'll actually need to program.
A \textbf{microcontroller} or \textbf{MCU} (\textit{Micro Controller Unit}) integrates onto a single chip what a
desktop splits across a whole motherboard:

\begin{itemize}
	\item
	      \textbf{CPU core} --- executes instructions (ours: ARM Cortex-M3);
	\item
	      \textbf{RAM} --- volatile working memory, tens of KB;
	\item
	      \textbf{Flash} --- non-volatile program storage; code runs directly
	      from it;
	\item
	      \textbf{Peripherals} --- timers, UART, GPIO, ADC, all on-chip;
	\item
	      One chip, one clock, one address space --- no external memory bus to
	      wire up, no BIOS, no disk.
\end{itemize}

Compared to desktop or handled general-purpose CPUs, we note that:

\begin{itemize}
	\item
	      A desktop CPU is \emph{just} the core: RAM, storage, and I/O are
	      separate chips reached over buses; an MMU gives every process its own
	      virtual address space;
	\item
	      An MCU has \textbf{one flat physical address map} shared by code,
	      data, and every peripheral register --- no MMU, no virtual memory (a
	      Cortex-M has at most an optional MPU);
	\item
	      On an MCU, ``the hardware'' is \emph{memory you write to}: there is no
	      driver layer between your code and a device register.
\end{itemize}

\subsubsection{Reset routine}

As an example, let's look into what happenes to a Cortex-M CPU at reset time.
After reset, the Cortex-M does not discover an operating system. It
reads two words from the vector table in Flash:

\begin{enumerate}
	\item
	      The initial value for the main stack pointer
	\item
	      The address of \texttt{Reset\_Handler}.
\end{enumerate}

The reset handler is our first code. It establishes the C runtime
environment before calling \texttt{main()}.

\subsubsection{Registers}

The CPU operates on values held in registers. Memory is accessed
explicitly with load and store instructions.

\begin{table}[h!]
	\center \rowcolors{2}{white}{black!10}
	\begin{tabular} { c | c }
		\bfseries Register role                  & \bfseries Meaning                           \\
		\hline
		Program Counter                          & Selects the next instruction                \\
		Stack Pointer                            & Locates calls, local variables, saved state \\
		General Registers / Programmer Registers & Hold operands and addresses                 \\
		Status Register                          & Records conditions and interrupt state      \\
	\end{tabular}
\end{table}

When an exception occurs, the processor automatically saves essential
state on the current stack so execution can later resume.

\subsubsection{Memory}

We can talk about the \textit{memory} available to an MCU.
Usually there are two kinds of memory:
\begin{itemize}
	\item \textbf{SRAM} / \textbf{DRAM}: not persistent, but extremely fast. Used as main program memory;
	\item \textbf{Flash} memory: read-only, but still very fast and fit (in embedded systems) for containing code. This is in accord with the usual definition of \textit{Harvard Architecture}. We should note that most general purpose devices use the opposite architecture, that is \textit{Von Neumann Architecture} (which use program memory both for code and data).
\end{itemize}

Flash keeps its contents without power. SRAM is fast writable working
memory whose contents disappear on reset or power loss.

\begin{verbatim}
Flash: vector table | .text | .rodata   | initial bytes for .data
SRAM:  .data        | .bss  | heap area | stack
\end{verbatim}

This split explains why startup code must copy some bytes and clear
others before ordinary C code can safely run.

\end{document}
